\documentclass[../../main.tex]{subfiles}
\begin{document}
\section{保守余向量场}\label{section:ch11-保守余向量场}

设 $M$ 为光滑流形。$M$ 上的一个余向量场 $\omega$ 称为\textbf{保守的}，若对 $M$ 中每个分段光滑闭曲线段 $\gamma$ 有
\begin{align*}
\int_\gamma\omega=0.
\end{align*}
它称为\textbf{恰当的}，若存在光滑函数 $f\in C^\infty(M)$ 使得 $\omega=df$。

由定理~\ref{theorem:ch11-恰当余向量场线积分}，每个恰当余向量场都是保守的。反过来如何？

\begin{proposition}\label{proposition:ch11-恰当蕴含保守}
每个恰当余向量场都是保守的。
\end{proposition}

\begin{proof}
若 $\omega=df$，则对任意闭曲线段 $\gamma$（端点同为 $p$），有 $\int_\gamma df=f(p)-f(p)=0$。
\end{proof}

下面的定理给出了保守余向量场的若干等价刻画。

\begin{theorem}\label{theorem:ch11-保守等价刻画}
设 $M$ 为带边或不带边的连通光滑流形，$\omega$ 为 $M$ 上的余向量场。下列条件等价：
\begin{enumerate}[(a)]
\item $\omega$ 是保守的；
\item $\omega$ 的线积分只依赖于曲线的端点（路径无关）；
\item $\omega$ 是恰当的。
\end{enumerate}
\end{theorem}

\begin{proof}
\eqref{align:ch11-恰当蕴含保守}~已证 (c) $\Rightarrow$ (a)。下证 (a) $\Rightarrow$ (b)：若 $\gamma_1$ 与 $\gamma_2$ 是从 $p$ 到 $q$ 的两条分段光滑曲线段，则先沿 $\gamma_1$ 从 $p$ 到 $q$、再沿反向 $\gamma_2$ 从 $q$ 回到 $p$，构成闭曲线。由保守性，
\begin{align*}
0=\int_{\gamma_1-\gamma_2}\omega=\int_{\gamma_1}\omega-\int_{\gamma_2}\omega,
\end{align*}
故 $\int_{\gamma_1}\omega=\int_{\gamma_2}\omega$。

下证 (b) $\Rightarrow$ (c)：假设线积分路径无关。固定任一基点 $p_0\in M$，定义函数 $f\colon M\to\mathbb{R}$ 为
\begin{align}
f(q)=\int_{p_0}^q\omega.\label{align:ch11-势函数}
\end{align}
我们需要证明 $f$ 光滑且 $df=\omega$。为此，令 $q_0\in M$ 任意，取以 $q_0$ 为中心的光滑坐标图 $(U,(x^i))$，在 $U$ 上把 $\omega$ 的坐标表示写为 $\omega=\omega_i\,dx^i$。我们需要证明
\begin{align}
\frac{\partial f}{\partial x^j}(q_0)=\omega_j(q_0)\label{align:ch11-势函数偏导}
\end{align}
对 $j=1,\dots,n$，这蕴含 $df_{q_0}=\omega_{q_0}$。

先设 $q_0\in\mathrm{Int}\,M$。固定 $j$，令 $\gamma\colon[-\varepsilon,\varepsilon]\to U$ 为在坐标下由 $\gamma(t)=(0,\dots,t,\dots,0)$（$t$ 在第 $j$ 位）定义的光滑曲线段，并取 $\varepsilon$ 足够小使 $\gamma[-\varepsilon,\varepsilon]\subseteq U$（\cref{figure:ch11-保守即恰当}）。令 $p_1=\gamma(-\varepsilon)$，定义新函数 $\widetilde f\colon M\to\mathbb{R}$ 为 $\widetilde f(q)=\int_{p_1}^q\omega$。注意\eqref{align:ch11-势函数}~蕴含对所有 $q\in M$，
\begin{align*}
f(q)-\widetilde f(q)=\int_{p_0}^q\omega-\int_{p_1}^q\omega=\int_{p_0}^q\omega+\int_q^{p_1}\omega=\int_{p_0}^{p_1}\omega,
\end{align*}
它不依赖于 $q$。于是 $\widetilde f$ 与 $f$ 只差一个常数，故只需证 $\partial\widetilde f/\partial x^j(q_0)=\omega_j(q_0)$。

\begin{figure}[htbp]
\centering
\includegraphics[width=0.6\textwidth]{fig-11-5.png}
\caption{保守余向量场即恰当的证明}\label{figure:ch11-保守即恰当}
\end{figure}

由构造 $\gamma'(t)=\partial/\partial x^j|_{\gamma(t)}$，于是
\begin{align*}
\omega_{\gamma(t)}(\gamma'(t))=\omega_i(\gamma(t))dx^i\left(\frac{\partial}{\partial x^j}\bigg|_{\gamma(t)}\right)=\omega_j(\gamma(t)).
\end{align*}
由于 $\gamma$ 限制到 $[-\varepsilon,t]$ 是从 $p_1$ 到 $\gamma(t)$ 的光滑曲线，有
\begin{align*}
\widetilde f\circ\gamma(t)=\int_{p_1}^{\gamma(t)}\omega=\int_{-\varepsilon}^t\omega_{\gamma(s)}(\gamma'(s))ds=\int_{-\varepsilon}^t\omega_j(\gamma(s))ds.
\end{align*}
于是由微积分基本定理，
\begin{align*}
\frac{\partial\widetilde f}{\partial x^j}(q_0)=\gamma'(0)\widetilde f=\frac{d}{dt}\bigg|_{t=0}\widetilde f\circ\gamma(t)
=\frac{d}{dt}\bigg|_{t=0}\int_{-\varepsilon}^t\omega_j(\gamma(s))ds=\omega_j(\gamma(0))=\omega_j(q_0).
\end{align*}
这表明 $df_{q_0}=\omega_{q_0}$ 当 $q_0\in\mathrm{Int}\,M$ 时。

若 $q_0\in\partial M$，则类似论证在半坐标下进行，结论相同。
\end{proof}

下面引入另一类余向量场。若在坐标图 $(U,(x^i))$ 中余向量场 $\omega=\omega_i\,dx^i$ 满足
\begin{align}
\frac{\partial\omega_i}{\partial x^j}=\frac{\partial\omega_j}{\partial x^i}\quad\text{对所有 }i,j,\label{align:ch11-闭条件}
\end{align}
则称 $\omega$ 在 $U$ 上\textbf{闭}。一个余向量场在 $M$ 上闭，若它在每个坐标图中闭。这是与坐标无关的条件（习题~\ref{exercise:ch11-40}）。

\textbf{例 11.43}。在穿孔平面 $\mathbb{R}^2\smallsetminus\{0\}$ 上，考虑余向量场
\begin{align*}
\omega=\frac{-y\,dx+x\,dy}{x^2+y^2}.
\end{align*}
直接计算表明它满足闭条件。它绕单位圆的线积分是 $2\pi$，故它不保守，因而不恰当。这说明闭余向量场不一定恰当。$\blacktriangleright$

\begin{proposition}\label{proposition:ch11-恰当蕴含闭}
每个恰当余向量场都是闭的。
\end{proposition}

\begin{proof}
设 $\omega=df$。在任意坐标图 $(U,(x^i))$ 中，$\omega_i=\partial f/\partial x^i$。由于偏导可交换次序，
\begin{align*}
\frac{\partial\omega_i}{\partial x^j}=\frac{\partial^2 f}{\partial x^j\partial x^i}=\frac{\partial^2 f}{\partial x^i\partial x^j}=\frac{\partial\omega_j}{\partial x^i},
\end{align*}
即 $\omega$ 闭。
\end{proof}

直接从定义检验一个余向量场是否闭，需在每个坐标图中验证\eqref{align:ch11-闭条件}。下面的命题给出了闭余向量场的一个与坐标无关的刻画，并附带表明只需在每点附近的某个坐标图中检验\eqref{align:ch11-闭条件}。

\begin{proposition}\label{proposition:ch11-闭的等价刻画}
设 $\omega$ 为带边或不带边光滑流形 $M$ 上的光滑余向量场。下列条件等价：
\begin{enumerate}[(a)]
\item $\omega$ 是闭的；
\item $\omega$ 在每点附近的某个光滑坐标图中满足\eqref{align:ch11-闭条件}；
\item 对任意开集 $U\subseteq M$ 与光滑向量场 $X,Y\in\mathfrak X(U)$，有
\begin{align}
X\big(\omega(Y)\big)-Y\big(\omega(X)\big)=\omega([X,Y]).\label{align:ch11-闭坐标无关}
\end{align}
\end{enumerate}
\end{proposition}

\begin{proof}
我们证 (a) $\Rightarrow$ (b) $\Rightarrow$ (c) $\Rightarrow$ (a)。(a) $\Rightarrow$ (b) 由闭余向量场定义立即得到。

下证 (b) $\Rightarrow$ (c)：设 (b) 成立，$U\subseteq M$，$X,Y\in\mathfrak X(U)$。只需在 $U$ 每点附近验证\eqref{align:ch11-闭坐标无关}。在任意坐标图 $(V,(x^i))$（$V\subseteq U$）中，写 $\omega=\omega_i dx^i$，$X=X^j\partial/\partial x^j$，$Y=Y^k\partial/\partial x^k$。利用 $[X,Y]$ 的坐标公式，直接计算即得。

下证 (c) $\Rightarrow$ (a)：设 (c) 成立。取坐标向量场 $X=\partial/\partial x^i$，$Y=\partial/\partial x^j$，则 $[X,Y]=0$，于是
\begin{align*}
0=\omega([X,Y])=X(\omega_j)-Y(\omega_i)=\frac{\partial\omega_j}{\partial x^i}-\frac{\partial\omega_i}{\partial x^j},
\end{align*}
即 $\omega$ 在该坐标图中闭。
\end{proof}
在单连通区域上，闭余向量场总是恰当的。我们先证明一个局部结果。

\begin{theorem}\label{theorem:ch11-星形区域闭即恰当}
设 $U\subseteq\mathbb{R}^n$ 是以 $c$ 为星形中心的开星形子集，即对每个 $x\in U$，从 $c$ 到 $x$ 的线段含于 $U$。若 $\omega$ 是 $U$ 上的光滑闭余向量场，则 $\omega$ 在 $U$ 上是恰当的。
\end{theorem}

\begin{proof}
对每个 $x\in U$，令 $\gamma_x\colon[0,1]\to U$ 为从 $c$ 到 $x$ 的线段：$\gamma_x(t)=c+t(x-c)$。定义
\begin{align}
f(x)=\int_{\gamma_x}\omega.\label{align:ch11-势函数星形}
\end{align}
我们需要证明 $f$ 是 $\omega$ 的势，即 $\partial f/\partial x^j=\omega_j$。先计算
\begin{align}
f(x)=\int_0^1\omega_{\gamma_x(t)}(\gamma_x'(t))dt=\int_0^1\omega_i(tx)x^i\,dt.\label{align:ch11-势函数星形积分}
\end{align}
（求和约定生效。）为计算 $f$ 的偏导数，注意被积函数在所有变量上光滑，故可在积分号下求导，得
\begin{align*}
\frac{\partial f}{\partial x^j}(x)=\int_0^1\left(t\frac{\partial\omega_i}{\partial x^j}(tx)x^i+\omega_j(tx)\right)dt.
\end{align*}
由于 $\omega$ 闭，这化为
\begin{align*}
\frac{\partial f}{\partial x^j}(x)&=\int_0^1\left(t\frac{\partial\omega_j}{\partial x^i}(tx)x^i+\omega_j(tx)\right)dt\\
&=\int_0^1\frac{d}{dt}\big(t\,\omega_j(tx)\big)dt=\big[t\,\omega_j(tx)\big]_{t=0}^{t=1}=\omega_j(x).
\end{align*}
\end{proof}

\begin{figure}[htbp]
\centering
\includegraphics[width=0.35\textwidth]{fig-11-6.png}
\caption{$\mathbb{R}^2$ 的一个星形子集}\label{figure:ch11-星形子集}
\end{figure}

\begin{corollary}[闭余向量场的局部恰当性]\label{corollary:ch11-局部恰当}
设 $\omega$ 为带边或不带边光滑流形 $M$ 上的闭余向量场。则 $M$ 的每一点都有一个邻域，在其上 $\omega$ 是恰当的。
\end{corollary}

\begin{proof}
任取 $p\in M$。假设意味着 $\omega$ 在某个包含 $p$ 的光滑坐标球或半球 $(U,\varphi)$ 中满足\eqref{align:ch11-闭条件}。由于球与半球是凸的，可对 $\omega$ 的坐标表示应用定理~\ref{theorem:ch11-星形区域闭即恰当}，得出存在函数 $f\in C^\infty(U)$ 使得 $\omega|_U=df$。
\end{proof}

定理~\ref{theorem:ch11-星形区域闭即恰当}~中构造势函数的关键是：我们能从 $c$ 出发沿一条随 $x$ 光滑变化的确定路径到达每一点 $x\in M$。在穿孔平面上的闭余向量场 $\omega$（例~11.43）情形中，正因有洞，这一点失败了：不可能选取一族光滑变化的路径，从固定基点出发，恰好到达定义域的每一点一次。在第~16~章中，我们将推广定理~\ref{theorem:ch11-星形区域闭即恰当}，证明每个闭余向量场在任意单连通流形上都恰当。

当你确实需要为已知恰当的余向量场计算势函数时，有一个简单得多的步骤几乎总是奏效。
\end{document}
